\section{Background: Maintenance Strategies}
\label{sec:background}

Maintenance strategies differ in what triggers an intervention and in the information they require (Table~\ref{tab:strategies}, Appendix~\ref{app:figures}). \textbf{Corrective} maintenance restores equipment after a failure; it needs no planning but exposes the plant to unplanned downtime and safety risks. \textbf{Preventive} maintenance acts at fixed time or usage intervals, which is easy to plan but ignores the state of each machine, so some components are replaced too early and others fail first. \textbf{Condition-based} maintenance (CBM) intervenes when a measured indicator, such as vibration or temperature, crosses a threshold. \textbf{Predictive} maintenance goes further and estimates the remaining useful life (RUL) with data-driven, physics-based or hybrid models. A RUL estimate, however, does not say which of several pending interventions should come first, nor how to weigh failure likelihood against safety, production or cost; many retrieved studies stopped at this stage and were excluded under E4. \textbf{Prescriptive} maintenance closes this gap by recommending the action, its timing and its resources, combining predictions with costs, constraints and expert knowledge. This is the level at which the present work is placed.

Two management frameworks are often combined with these strategies. Reli\-a\-bil\-i\-ty-Centred Maintenance (RCM) defines the policy of each asset through an analysis of its functions, failure modes and consequences; its tools, FMEA and criticality analysis, introduced the ranking of assets by severity and likelihood that underlies many of the models in Section~\ref{sec:criteria-methods}. Total Productive Maintenance (TPM) involves operators in equipment care and measures performance through Overall Equipment Effectiveness (OEE), bringing the production cost of downtime into maintenance decisions.